\section*{الفصل \ref{ch:b3:gfd-foundations} --- ديناميكا الموائع الجيوفيزيائية}

\begin{solution}{exo:b3:gfd-foundations:1}
(أ)~$f = 2\Omega\sin\varphi$: \num{7.292e-5}، \num{1.031e-4}
و\qty{1.263e-4}{s^{-1}}؛ $2\pi/f$: 23.93، 16.92 و\qty{13.82}{h}.
(ب)~$U/f$: 2.74، 1.94 و\qty{1.58}{km}. (ج)~عند \ang{30}، $\sin\varphi =
1/2$ و$2\pi/f = 2\pi/\Omega$: الدور القصوري هو اليوم النجمي،
\qty{23.93}{h}.
\end{solution}

\begin{solution}{exo:b3:gfd-foundations:2}
مع $f = \qty{1.031e-4}{s^{-1}}$: زوبعة الغبار \num{9.7e3}، نسيم البحر 2.4،
المنخفض 0.15، والدوّامة المحيطية \num{2.4e-4}. يقترب المنخفض والدوّامة من
الاتزان الجيوستروفي؛ أمّا زوبعة الغبار فتدور في أيّ الاتجاهين، غافلةً عن
دوران الأرض، ولا يشعر به نسيم البحر إلّا على مدى يوم.
\end{solution}

\begin{solution}{exo:b3:gfd-foundations:3}
$\partial u_g/\partial z = (g/fT)\,|\partial T/\partial y| = 9.81\times
\num{5e-6}/(\num{1.031e-4}\times250) = \qty{1.90e-3}{s^{-1}}$؛ على مدى
\qty{10}{km}، تزداد الريح الغربية بمقدار \qty{19}{m/s}.
\end{solution}

\begin{solution}{exo:b3:gfd-foundations:4}
$E = \num{50e-8}$: $N = \qty{2.21e-3}{s^{-1}}$، دور قدره \qty{47}{min}،
\num{1.27} دورةً في الساعة. $E = \num{1000e-8}$: $N = \qty{9.90e-3}{s^{-1}}$،
\qty{10.6}{min}، 5.7 دورةً في الساعة. $E = \num{e-8}$: $N =
\qty{3.13e-4}{s^{-1}}$، \qty{5.6}{h}، 0.18 دورةً في الساعة.
\end{solution}

\begin{solution}{exo:b3:gfd-foundations:5}
(أ)~$W \sim UH/L = \qty{0.1}{m/s}$ و$\mathrm Dw/\mathrm Dt \sim UW/L =
\qty{e-6}{m/s^2}$، \num{1.0e-7} من $g$. حدّ الضغط، $fUL/H =
\qty{0.103}{m/s^2}$، أكبر بمقدار \num{1e5}: النسبة \num{9.7e-6}، هي
$\mathrm{Ro}(H/L)^2$ مع $\mathrm{Ro} = 0.097$. (ب)~$\tilde fu =
2\Omega\cos\ang{45}\times10 = \qty{1.03e-3}{m/s^2}$، \num{1.05e-4} من $g$.
\end{solution}

\begin{solution}{exo:b3:gfd-foundations:6}
$fr = \qty{30}{m/s}$. (أ)~$\mathrm{Ro}_g = 0.667$، $V = 40/(1 + \sqrt{3.667})
= \qty{13.7}{m/s}$، أبطأ بمقدار الثلث من الحالة الجيوستروفية. (ب)~$\mathrm{Ro}_g =
0.167$، $V = 10/(1 + \sqrt{0.333}) = \qty{6.34}{m/s}$. (ج)~$V_g \le fr/4 =
\qty{7.5}{m/s}$، ثم $V = fr/2 = \qty{15}{m/s}$، أي ضعف
السرعة الجيوستروفية.
\end{solution}

\begin{solution}{exo:b3:gfd-foundations:7}
(أ)~$\Omega = 2\pi/6 = \qty{1.047}{rad/s}$، $f = \qty{2.094}{s^{-1}}$.
$\mathrm{Ro} = 0.01/(2.094\times0.05) = 0.095$ و$\mathrm{Ek} =
\num{e-6}/(2.094\times0.04) = \num{1.19e-5}$: كلاهما صغير، فيخضع الجريان
\omterm{thm:b3:gfd-foundations:taylor-proudman}{لمبرهنة تايلور--برودمان}، ويتكوّن \omterm{thm:b3:gfd-foundations:taylor-proudman}{عمود تايلور}. (ب)~$0.2/\sqrt{\num{e-6}
\times1.047} = \qty{195}{s}$، \qty{3.3}{min}، 33 دورةً من دوران الطاولة.
\end{solution}

\begin{solution}{exo:b3:gfd-foundations:8}
$\partial\bar T/\partial y = -f\,\Delta u/(R\ln(850/250))$، مع
$\ln(850/250) = 1.224$. عند \ang{30}~شمالًا، $-\num{7.29e-5}\times36.12/(286.9
\times1.224) = \qty{-7.50}{K}$ لكلّ \qty{1000}{km}؛ عند \ang{45}~شمالًا،
$-\num{1.031e-4}\times17.18/351.1 = \qty{-5.05}{K}$ لكلّ \qty{1000}{km}. القصّ
يتناسب مع التدرّج مقسومًا على $f$: التدرّج أكبر بمقدار 1.5
مرّة عند \ang{30}، و$f$ أصغر بمقدار 1.4 مرّة، فالقصّ أكبر بمقدار 2.1
مرّة.
\end{solution}

\begin{solution}{exo:b3:gfd-foundations:9}
(أ)~$L_R = 0.0113\times11\,000/\num{1.031e-4} = \qty{1206}{km}$؛
$0.00357\times1000/\num{9.47e-5} = \qty{37.7}{km}$، أصغر بمقدار 32 مرّة.
(ب)~$\beta = 2\Omega\cos\varphi/a$: \num{1.62e-11}
و\qty{1.74e-11}{m^{-1}.s^{-1}}. $\beta L_R/f = 0.19$ للغلاف الجوي، 0.007
وللمحيط.
\end{solution}

\begin{solution}{exo:b3:gfd-foundations:10}
(أ)~$\mathrm dw/\mathrm dt = -(r + \mathrm if)w$، فإنّ $w = U\mathrm
e^{-rt}\mathrm e^{-\mathrm ift}$: تدور السرعة بالمعدّل $f$ بينما
يتناقص قدرها وفق $\mathrm e^{-rt}$. ويكون الموضع $x + \mathrm iy =
U(1 - \mathrm e^{-(r + \mathrm if)t})/(r + \mathrm if)$ يتلوّى حلزونيًا نحو
النقطة $U/(r + \mathrm if)$. (ب)~$x + \mathrm iy = U(1 - \mathrm e^{-(r +
\mathrm if)t})/(r + \mathrm if) + u_0t$: يتمدّد الحلزون إلى حلقات.
(ج)~$0.2\,\mathrm e^{-1} = \qty{0.074}{m/s}$؛ انجرف المركز بمقدار $u_0t =
0.05\times4\times86\,400 = \qty{17.3}{km}$ شرقًا.
\end{solution}

\begin{solution}{exo:b3:gfd-foundations:11}
$f = \qty{9.37e-5}{s^{-1}}$؛ $\partial u/\partial z = (g/f\rho_0)\,\partial
\rho/\partial y = 9.81\times10^{-5}/(\num{9.37e-5}\times1025) =
\qty{1.02e-3}{s^{-1}}$. مع $u = 0$ عند \qty{1000}{m} يبلغ تيّار السطح
\qty{1.02}{m/s}، شرقًا: في نصف الكرة الشمالي يجري التيّار
مع الماء الخفيف على يمينه.
\end{solution}

\begin{solution}{exo:b3:gfd-foundations:12}
(أ)~$N^2 = (g/\rho)\,\Delta\sigma_\theta/\Delta z$. 0--\qty{100}{m}:
$9.81\times0.668/(1026.4\times100)$، $N = \qty{7.99e-3}{s^{-1}}$، 4.6
دورةً في الساعة؛ 100--\qty{500}{m}: \qty{3.02e-3}{s^{-1}}، 1.7 دورةً في
الساعة؛ 500--\qty{1000}{m}: \qty{3.24e-3}{s^{-1}}، 1.9 دورةً في الساعة.
(ب)~الطبقة العليا \qty{100}{m}: تحتوي على ماء السطح الدافئ فوق
الطبقة الانتقالية الكثافية الموسمية. بين 100 و\qty{500}{m} يكون الماء أكثر
انتظامًا من الأسفل. (ج)~$E = N^2/g$ هي \num{650e-8}، \num{93e-8}
و\qty{107e-8}{m^{-1}}، وكلّها ضمن مدى المحيط العلوي.
\end{solution}

\begin{solution}{pb:b3:gfd-foundations:1}
\textbf{1.} $f = 2\times\num{7.292115e-5}\sin\ang{48} =
\qty{1.084e-4}{s^{-1}}$؛ $2\pi/f = \qty{16.1}{h}$.

\textbf{2.} حوض الاستحمام \num{3.1e3}؛ الخزّان ($f = 4\pi/10 = \qty{1.257}{s^{-1}}$)
0.020؛ نسيم البحر 2.3؛ المنخفض 0.14؛ الدوّامة المحيطية \num{2.3e-4}.

\textbf{3.} لا يشعر حوض الاستحمام بشيء من دوران الأرض؛ ويشعر نسيم البحر
به إلّا بعد بضع ساعات. أمّا الخزّان والمنخفض والدوّامة المحيطية فلها
أعداد روسبي صغيرة، وهي قريبة من الاتزان الجيوستروفي.

\textbf{4.} الخزّان $\num{e-6}/(1.257\times0.04) = \num{2.0e-5}$؛ المنخفض
$10/(\num{1.084e-4}\times10^8) = \num{9.2e-4}$؛ الدوّامة المحيطية
$\num{e-3}/(\num{1.084e-4}\times10^6) = \num{9.2e-6}$.

\textbf{5.} $\delta/H = \sqrt{2\nu/(fH^2)} = \sqrt{2\,\mathrm{Ek}}$. الخزّان
$\sqrt{2\times\num{e-6}/1.257} = \qty{1.3}{mm}$؛ الهواء $\sqrt{20/\num{1.084e-4}}
= \qty{430}{m}$.

\textbf{6.} يؤثّر الاحتكاك في طبقات بسمك بضعة أجزاء في الألف من العمق، في
الخزّان والمحيط، وبضعة أجزاء في المئة في الغلاف الجوي؛ وفوقها يكون الجريان
شبه خالٍ من الاحتكاك.

\textbf{7.} $\rho = p/(RT) = 101\,290/(286.9\times288.1) =
\qty{1.225}{kg/m^3}$.

\textbf{8.} $u_g = (1/\rho f)\,|\partial p/\partial y| = (400/300\,000)/(1.225
\times\num{1.084e-4}) = \qty{10.0}{m/s}$، من الغرب، مع انخفاض الضغط
شمالًا على اليسار.

\textbf{9.} $\mathrm{Ro}_g = 10.04/(\num{1.084e-4}\times600\,000) = 0.154$;
$V = 2\times10.04/(1 + \sqrt{1.618}) = \qty{8.84}{m/s}$, \qty{12}{\percent}
أبطأ.

\textbf{10.} $V = 2\times10.04/(1 + \sqrt{0.382}) = \qty{12.4}{m/s}$. وأكبر
قيمة للمقدار $V_g$ هي $fr/4 = \qty{16.3}{m/s}$، وتدرّج الضغط $\rho fV_g =
\qty{2.16e-3}{Pa/m}$، \qty{6.5}{hPa} لكلّ \qty{300}{km}.

\textbf{11.} $g/c_p = \qty{9.77}{K/km}$. $N^2 = (9.81/255.7)\times(9.77 -
6.49)\times10^{-3} = \qty{1.26e-4}{s^{-2}}$، $N = \qty{0.0112}{s^{-1}}$، أي
زمن دوريّ قدره \qty{9.3}{min}.

\textbf{12.} عند $\Gamma = 0$، $N^2 = (g/T)(g/c_p) = g^2/(c_pT) =
\qty{4.43e-4}{s^{-2}}$، $N = \qty{0.0210}{s^{-1}}$، أي 1.9 ضعف
القيمة في التروبوسفير.

\textbf{13.} $L_R = 0.0112\times10^4/\num{1.084e-4} = \qty{1035}{km}$، وهو
يساوي حجم المنخفض.

\textbf{14.} $g' = \qty{0.0147}{m/s^2}$; $\sqrt{g'h} = \qty{0.54}{m/s}$;
$L_R = 0.542/\num{1.084e-4} = \qty{5.0}{km}$.

\textbf{15.} نعم: البحيرة، التي عرضها \qty{12}{km}، أكبر من ضعف
نصف قطر روسبي لها، لذلك تنحرف أمواجها الداخلية وتياراتها بفعل
قوّة كوريوليس.

\textbf{16.} $(249.4 - 253.0)/600 = \qty{-6.0}{K}$ لكلّ \qty{1000}{km}.

\textbf{17.} $\ln(1013/265) = 1.3409$. الجنوب:
$286.9\times253.0\times1.3409/9.81 = \qty{9922}{m}$؛ الشمال \qty{9781}{m}.
يقع سطح \qty{265}{hPa} أخفض بمقدار \qty{141}{m} في الشمال، بانحدار
\num{2.35e-4} بالنسبة إلى السطح السفلي.

\textbf{18.} $\Delta u = (g/f)\times\text{الانحدار} = 9.81\times\num{2.35e-4}/
\num{1.084e-4} = \qty{21.3}{m/s}$.

\textbf{19.} $(R/f)\ln(p_b/p_t)\,|\partial\bar T/\partial y| =
(286.9/\num{1.084e-4})\times1.3409\times\num{6.0e-6} = \qty{21.3}{m/s}$، وهي
نفسها: الحسابان هما حساب واحد.

\textbf{20.} \omterm{def:b3:gfd-foundations:baroclinic}{باروكليني}: تتغيّر درجة الحرارة، ومن ثمّ الكثافة، على طول
أسطح الضغط. فلو كان الجريان باروتروبيًّا لما تغيّرت الريح مع الارتفاع:
\qty{10.0}{m/s}.

\textbf{21.} $(g/f\bar T)\,|\partial T/\partial y|\,\Delta z = 9.81\times
\num{6.0e-6}\times10^4/(\num{1.084e-4}\times255.7) = \qty{21.2}{m/s}$,
في حدود \qty{0.3}{\percent}.

\textbf{22.} $10.04 + (286.9/\num{1.084e-4})\times\num{6.0e-6}\times
\ln(1013/541) = 10.04 + 9.96 = \qty{20.0}{m/s}$.

\textbf{23.} $\partial u/\partial z = 21.3/10^4 = \qty{2.13e-3}{s^{-1}}$;
$\mathrm{Ri} = \num{1.26e-4}/\num{4.54e-6} = 28$، أكبر بكثير من 0.25: القصّ
لا يجعل التروبوسفير الحرّ مضطربًا.

\textbf{24.} ريح أورين، \qty{31}{m/s}، أقوى من المتوسط النطاقي
عند \ang{45}~شمالًا وأضعف منه عند \ang{30}~شمالًا. تتذبذب نفّاثة الجبهة القطبية
شمالًا وجنوبًا؛ وفي أيّ لحظة يتركّز تدرّجها وريحها
على نطاق ضيّق، بينما يؤدّي أخذ المتوسط حول دائرة العرض إلى توزيعهما.

\textbf{25.} $10.0 + 21.3$: الريح الجيوستروفية عند \qty{10}{km} فوق
السهوب هي \textbf{\qty[text-series-to-math=true]{31}{m/s}}، من الغرب
(\qty{31.3}{m/s} حسب الحساب).
\end{solution}
